% Estilo visual do PPC — pacotes, cores e o tema de tabela "ppc" usado por
% ppcgen.geradores.tabelas/fluxo/representacao_grafica.
%
% Adaptado do template original (include/settings/Settings.tex, criado por
% David Frisk/2018, modificado por Igor S Peretta/2023) removendo apenas as
% macros de identificação do curso (\nomecursonovo, \ppcversao), que agora
% são geradas a partir de config/curso.yaml em latex/gerado/curso_macros.tex
% — o restante do estilo visual é genérico e não específico de nenhum
% curso, por isso foi reaproveitado em vez de reescrito.

\usepackage{xspace}

% BASIC SETTINGS
\usepackage{moreverb}
\usepackage{textcomp}
\usepackage{lmodern}
\usepackage{helvet}
\usepackage[T1]{fontenc}
\usepackage[brazil]{babel}
\usepackage[utf8]{inputenc}
\usepackage{amsmath}
\usepackage{amssymb}
\usepackage{graphicx}
\usepackage{subfig}
\numberwithin{equation}{chapter}
\numberwithin{figure}{chapter}
\numberwithin{table}{chapter}
\usepackage{listings}
\usepackage{multirow}
\usepackage[table]{xcolor}
\usepackage[top=3cm, bottom=3cm, inner=3cm, outer=3cm]{geometry}
\usepackage{eso-pic}
\newcommand{\backgroundpic}[3]{
    \put(#1,#2){
    \parbox[b][\paperheight]{\paperwidth}{
    \centering
    \includegraphics[width=\paperwidth,height=\paperheight,keepaspectratio]{#3}}}}
\usepackage{float}
\usepackage{parskip}
\usepackage{wrapfig}
\usepackage{hyperref}
\usepackage{tikz}
\usetikzlibrary{fit}
\usepackage{enumerate}
\usepackage{csquotes}
\usepackage{soul}
\usepackage{caption}
\usepackage{makecell}

% CORES
\definecolor{AzulClaro}{RGB}{41,102,189}
\definecolor{AzulEscuro}{RGB}{14,64,151}
\definecolor{CinzaClaro}{RGB}{232,231,231}
\definecolor{CinzaMeio}{RGB}{172,175,180}
\definecolor{CinzaEscuro}{RGB}{113,119,130}
\definecolor{Verde}{RGB}{1,66,10}
\definecolor{Amarelo}{RGB}{232,171,2}
\definecolor{Laranja}{RGB}{204,85,0}
\definecolor{AzulMarca}{RGB}{220,230,241}
\definecolor{AmareloVivo}{RGB}{255,255,0}

\hypersetup{
    colorlinks=true,
    linkcolor=AzulClaro,
    filecolor=AzulClaro,
    citecolor=AzulEscuro,
    urlcolor=AzulClaro,
}

\usepackage[explicit]{titlesec}
\titleformat{\chapter}[display]
  {\normalfont}{}{0pt}{%
    \begin{tikzpicture}
    \node[fill=AzulEscuro, minimum width=1.1cm, minimum height=1.1cm, rounded corners, anchor=west,
      font=\color{white}\fontsize{18}{24}\selectfont\bfseries](chapnumber){\thechapter};
    \node[anchor=west, text width=\textwidth-4cm, font=\color{AzulEscuro}\bfseries\Large]
      at ([xshift=10pt]chapnumber.east){#1};
    \fill[overlay, draw=none, line width=0pt, rounded corners=0pt,
      left color=AzulEscuro, right color=AzulClaro!25]
      ([yshift=12pt]chapnumber.north west) rectangle ++(\textwidth,-2pt);
    \end{tikzpicture}%
  }
\titleformat{name=\chapter,numberless}[display]
  {\normalfont}{}{0pt}{%
    \begin{tikzpicture}
    \node[anchor=west, inner sep=0pt, outer sep=0pt, text width=\textwidth,
      font=\color{AzulEscuro}\bfseries\Large](chaptitle){#1};
    \fill[overlay, draw=none, line width=0pt, rounded corners=1pt,
      left color=AzulEscuro, right color=AzulClaro!25]
      ([yshift=-15pt]chaptitle.south west) rectangle ++(\textwidth,-2pt);
    \end{tikzpicture}%
  }
\titleformat{\section}[display]
  {\normalfont}{}{0pt}{%
    \begin{tikzpicture}
    \node[fill=AzulClaro, minimum width=0.9cm, minimum height=0.9cm, rounded corners, anchor=west,
      font=\color{white}\fontsize{12}{14}\selectfont\bfseries](secnumber){\thesection};
    \node[anchor=west, text width=\textwidth-4cm, font=\color{AzulClaro}\bfseries\large]
      at ([xshift=8pt]secnumber.east){#1};
    \end{tikzpicture}%
  }
\titleformat{name=\section,numberless}[display]
  {\normalfont}{}{0pt}{%
    \begin{tikzpicture}
    \node[anchor=west, inner sep=0pt, outer sep=0pt, text width=\textwidth,
      font=\color{AzulClaro}\bfseries\normalsize](sectitle){#1};
    \end{tikzpicture}%
  }

\usepackage{tabularray}
\NewTblrTheme{ppc}{%
    \DefTblrTemplate{caption-tag}{normal}{\textbf{\textcolor{AzulEscuro}{Quadro\hspace{0.25em}\thetable}}}
    \DefTblrTemplate{caption-sep}{normal}{\textbf{\textcolor{AzulEscuro}{:}}\enskip}
    \DefTblrTemplate{caption-text}{normal}{\InsertTblrText{caption}}
    \DefTblrTemplate{contfoot-text}{normal}{\tiny \textcolor{AzulClaro}{Continua na próxima página}}
    \DefTblrTemplate{conthead-text}{normal}{\tiny \textcolor{AzulClaro}{(Continuação)}}
    \SetTblrTemplate{caption-tag}{normal}
    \SetTblrTemplate{caption-sep}{normal}
    \SetTblrTemplate{caption-text}{normal}
    \SetTblrTemplate{contfoot-text}{normal}
    \SetTblrTemplate{conthead-text}{normal}
}

\usepackage[labelfont=bf, textfont=normal, justification=justified, singlelinecheck=false]{caption}

\setcounter{tocdepth}{5}
\setcounter{secnumdepth}{5}

\usepackage{fancyhdr}
\pagestyle{fancy}
\renewcommand{\chaptermark}[1]{\markboth{\thechapter.\space#1}{}}
\fancyhf{}
\fancyhead[LE,RO]{\nouppercase{\leftmark}}
\fancyfoot[LE,RO]{\thepage}
\fancypagestyle{plain}{
  \fancyhf{}
  \renewcommand{\headrulewidth}{0pt}
  \fancyfoot[LE,RO]{\thepage}
}
\setlength{\headheight}{15pt}

\usepackage{footnotehyper}
\makesavenoteenv{longtblr}
\usepackage[inline]{enumitem}
\newcommand{\listspace}{\setlist[enumerate]{itemsep=5pt,parsep=0pt,topsep=0pt,partopsep=0pt}}

\usepackage{lscape}
\usepackage{xifthen}
\newcommand\nohyph{\hyphenpenalty=10000\relax\exhyphenpenalty=10000\relax}

\usepackage{etoolbox}
\AtBeginDocument{%
    \renewcommand{\familydefault}{\sfdefault}
    \normalfont
}

% Macros do diagrama de fluxo (representação gráfica) — genéricas,
% reaproveitadas do template original.
\newcommand{\ccheader}[6]{
    \begin{tikzpicture}
        \draw (0,0.5) rectangle node[font=\bfseries\sffamily\scriptsize, text width=3cm, align=center]{#1Período} (3,1);
        \draw (0,0) rectangle node[font=\bfseries\sffamily\tiny, text width=0.75cm, align=center]{CHT} (0.6,0.5);
        \draw (0.6,0) rectangle node[font=\bfseries\sffamily\tiny, text width=0.6cm, align=center]{CHP} (1.2,0.5);
        \draw (1.2,0) rectangle node[font=\bfseries\sffamily\tiny, text width=0.6cm, align=center]{CHD} (1.8,0.5);
        \draw (1.8,0) rectangle node[font=\bfseries\sffamily\tiny, text width=0.6cm, align=center]{CHE} (2.4,0.5);
        \draw (2.4,0) rectangle node[font=\bfseries\sffamily\tiny, text width=0.6cm, align=center]{TOT} (3,0.5);
        \draw (0,-0.5) rectangle node[font=\bfseries\sffamily\scriptsize, text width=0.6cm, align=center]{#2} (0.6,0);
        \draw (0.6,-0.5) rectangle node[font=\bfseries\sffamily\scriptsize, text width=0.6cm, align=center]{#3} (1.2,0);
        \draw (1.2,-0.5) rectangle node[font=\bfseries\sffamily\scriptsize, text width=0.6cm, align=center]{#4} (1.8,0);
        \draw (1.8,-0.5) rectangle node[font=\bfseries\sffamily\scriptsize, text width=0.6cm, align=center]{#5} (2.4,0);
        \draw (2.4,-0.5) rectangle node[font=\bfseries\sffamily\scriptsize, text width=0.6cm, align=center]{#6} (3,0);
    \end{tikzpicture}
}

\newcommand{\ccbloco}[9]{
    \begin{tikzpicture}
        \draw (0,0.5) rectangle node[font=\sffamily\scriptsize, yshift=-2pt, text width=2.8cm, align=center]{#2} (3,2.5);
        \draw (1.5,2.5) node[font=\bfseries\sffamily\scriptsize, yshift=-7pt, anchor=mid]{#1};
        \draw (0,0) rectangle node[font=\sffamily\scriptsize, text width=1cm, align=center]{#3} (0.6,0.5);
        \draw (0.6,0) rectangle node[font=\sffamily\scriptsize, text width=1cm, align=center]{#4} (1.2,0.5);
        \draw (1.2,0) rectangle node[font=\sffamily\scriptsize, text width=1cm, align=center]{#5} (1.8,0.5);
        \draw (1.8,0) rectangle node[font=\sffamily\scriptsize, text width=1cm, align=center]{#6} (2.4,0.5);
        \draw (2.4,0) rectangle node[font=\sffamily\scriptsize, text width=1cm, align=center]{#7} (3,0.5);
        \ifthenelse{\isempty{#8}}
        {\draw (-0.475,1.75) node[font=\sffamily\scriptsize, yshift=-7pt, anchor=center, align=right, text width=2ex]{};}
        {\draw (-0.475,1.75) node[font=\sffamily\scriptsize, yshift=-7pt, anchor=center, align=right, text width=2ex]{\textbf{#8}};
         \draw (-0.10,1.75) node[font=\sffamily\footnotesize, yshift=-7pt, anchor=mid, rotate=90]{$\boldsymbol{\blacktriangledown}$};}
        \ifthenelse{\isempty{#9}}
        {\draw (-0.475,0.5) node[font=\sffamily\scriptsize, yshift=-7pt, anchor=center, align=right, text width=2ex]{};}
        {\draw (-0.475,0.5) node[font=\sffamily\scriptsize, yshift=-7pt, anchor=center, align=right, text width=2ex]{\textbf{#9}};
         \draw (-0.30,0.5) node[font=\sffamily\footnotesize, yshift=-7pt, anchor=mid, rotate=90]{\rotatebox{180}{$\boldsymbol{\bigstar}$}};}
    \end{tikzpicture}
}
